\chapter{Barrières et digitales}\label{ch:dv:barriers-and-digitals}

Un desk a vendu une \omterm{def:dv:barriers-and-digitals:digital}{option digitale} qui verse 10 millions si l'indice clôture au-dessus de 5\,000. Une minute avant
la clôture, l'indice vaut 4\,999.9. L'option vaut 4.9 millions, et un mouvement d'un point d'indice dans un sens
ou dans l'autre la modifie de 1.25 million. Pour la couvrir, le desk aurait besoin d'une position sur l'indice
équivalente à 6.25 milliards, achetée ou vendue durant la dernière minute de séance puis défaite si l'indice
repart en arrière. Aucun desk ne fait cela. La discontinuité de la digitale la rend non couvrable au strike à
l'approche de l'échéance. La même discontinuité se retrouve dans toute option à barrière, que le marché négocie en gros
volumes sur les devises et au sein des produits structurés. Ce chapitre évalue les digitales et les barrières et explique
comment les desks couvrent ce qui ne peut pas être couvert en delta~: par une surcouverture, par des portefeuilles
statiques de vanilles, et par une barrière décalée pour un suivi discret. Il se termine sur le risque qu'aucune de ces méthodes ne supprime~: un prix qui saute par-dessus la barrière.

\section{Digitales et leur couverture}

\begin{definition}[Option digitale]\label{def:dv:barriers-and-digitals:digital}
Une \emph{option digitale}\index{option digitale} (binary option) verse un montant fixe si le sous-jacent termine au-dessus
(call) ou au-dessous (put) du strike, et rien sinon. Sous Black--Scholes, un call digital au comptant payant 1 a pour prix
$e^{-rT}\Phi(d_2)$.
\end{definition}

Une digitale est la limite de call spreads, $\bigl(C(K-\varepsilon)-C(K)\bigr)/\varepsilon$ lorsque $\varepsilon\to0$,
si bien que son prix est l'opposé de la dérivée du prix du call par rapport au strike. En présence d'un smile, la volatilité du call dépend du
strike, et la dérivation donne
\[
D(K)=-\frac{\partial C}{\partial K}=D_{\mathrm{BS}}(K,\sigma(K))-\mathcal V(K)\,\frac{\partial\sigma}{\partial K}.
\]
Avec un skew baissier $\partial\sigma/\partial K<0$, le call digital vaut plus que son prix de Black--Scholes à la
volatilité du strike. Le terme de skew constitue souvent la part la plus importante de l'écart entre les cotations des desks.

\begin{example}[Une digitale à trois mois sur la surface du chapitre~9]\label{ex:dv:barriers-and-digitals:smile}
Ramenons la surface du chapitre~9 à un indice à 5\,000. Un call digital à trois mois, à la monnaie, sur 10 millions vaut
4.84 millions avec Black--Scholes à la volatilité à la monnaie de 16.4~\%. Avec le smile, il vaut 5.77 millions~:
le terme de \omterm{def:dv:greeks-and-the-hedging-pnl:greeks}{vega} ajoute 0.93 million, soit près du cinquième du prix.
\end{example}

À l'approche de l'échéance, le delta de la digitale, $e^{-rT}\varphi(d_2)/(S\sigma\sqrt T)$ par unité payée, croît
comme $1/\sqrt T$ au strike et s'annule en s'en éloignant (\cref{fig:dv:barriers-and-digitals:digital}, à gauche).
Son gamma change de signe au strike. Un desk vendeur d'une digitale proche du strike est soit vendeur, soit acheteur
d'une convexité considérable, selon le côté du strike où se trouve l'indice, et il ne peut pas trader assez vite
pour suivre. La réponse pratique consiste à remplacer la digitale par un call spread.

\begin{definition}[Surcouverture par call spread]\label{def:dv:barriers-and-digitals:overhedge}
La \emph{surcouverture par call spread}\index{surcouverture par call spread} évalue et couvre un call digital vendu
comme le call spread qui le domine~: $N/\varepsilon$ calls de strike $K-\varepsilon$ moins autant de calls de strike $K$. Le spread
verse au moins autant que la digitale en tout point, et l'écart de prix, la surcouverture, est facturé à l'acheteur.
\end{definition}

La largeur $\varepsilon$ arbitre entre prix et risque. Un spread étroit est bon marché, mais il conserve l'essentiel de la convexité de la digitale~: son delta peut atteindre
$N/\varepsilon$ par point d'indice. Un spread large est facile à couvrir, mais il fait payer au client un versement qu'il n'a pas acheté.
Un desk fixe la largeur à partir de la perte maximale qu'il accepte par point du cours de règlement-livraison, soit le point
où la pente du spread devient le risque. Sur l'exemple, des largeurs de 50, 100, 200
et 400 points d'indice coûtent 0.25, 0.48, 0.89 et 1.56 million en sus des 5.77 millions de la digitale.

\begin{omfigure}
\begin{tikzpicture}
\begin{groupplot}[group style={group size=2 by 1,horizontal sep=1.6cm},
  width=6.6cm,height=5cm,
  tick label style={font=\small},label style={font=\small},grid=major,grid style={gray!25},table/col sep=comma]
\nextgroupplot[xlabel={indice},ylabel={delta (millions par point)},xmin=4900,xmax=5100,ymin=0,ymax=0.18,
  xtick={4900,5000,5100},yticklabel style={/pgf/number format/fixed,/pgf/number format/precision=2},
  title={\small delta d'une digitale de 10 millions},
  legend columns=3,legend style={font=\footnotesize,at={(0.5,-0.32)},anchor=north,draw=none,fill=none}]
\addplot[omProp,thick,dashed] table[x=s,y=week]{figdata/derivatives/15-barriers-and-digitals/digital_delta.csv};
\addplot[omDef,thick] table[x=s,y=day]{figdata/derivatives/15-barriers-and-digitals/digital_delta.csv};
\addplot[omThm,very thick] table[x=s,y=hour]{figdata/derivatives/15-barriers-and-digitals/digital_delta.csv};
\legend{1 semaine,1 jour,1 heure}
\nextgroupplot[xlabel={indice à l'échéance},ylabel={paiement (millions)},xmin=4600,xmax=5200,ymin=-0.5,ymax=11,
  xtick={4600,4800,5000,5200},
  title={\small digitale et surcouverture},
  legend columns=1,legend cell align=left,legend style={font=\footnotesize,at={(0.5,-0.32)},anchor=north,draw=none,fill=none}]
\addplot[omThm,very thick] table[x=s,y=digital]{figdata/derivatives/15-barriers-and-digitals/overhedge.csv};
\addplot[omDef,thick,dashed] table[x=s,y=spread]{figdata/derivatives/15-barriers-and-digitals/overhedge.csv};
\legend{digitale,{call spread 4\,800--5\,000}}
\end{groupplot}
\end{tikzpicture}

\omcaption{À gauche~: le delta d'une digitale de 10 millions de strike 5\,000 (volatilité de 20~\%) à mesure que
l'échéance approche~; une minute avant la clôture, il atteint 1.25 million par point, hors de cette échelle. À droite~: le paiement de la digitale et
le call spread qui la surcouvre~; le versement supplémentaire entre 4\,800 et 5\,000 est ce que facture la surcouverture. Données~: code du
chapitre.}\label{fig:dv:barriers-and-digitals:digital}
\end{omfigure}

\section{Barrières par réflexion}

Une option à barrière (livre~2, chapitre~19) s'active ou se désactive lorsque le sous-jacent touche un niveau $H$. Avec
un suivi continu et une dynamique de Black--Scholes, les huit barrières simples admettent des formes fermées, qui
découlent toutes du principe de réflexion (livre~4, chapitre~2).

\begin{proposition}[Le call down-and-out par méthode des images]\label{prop:dv:barriers-and-digitals:image}
Pour $K\ge H$, avec le portage $b=r-q$ et $\nu=b-\frac12\sigma^2$, le call down-and-out vaut
\[
C_{\mathrm{DO}}(S)=C_{\mathrm{BS}}(S,K)-\Bigl(\frac HS\Bigr)^{2\nu/\sigma^2}C_{\mathrm{BS}}\Bigl(\frac{H^2}S,K\Bigr).
\]
\end{proposition}
\begin{proof}
La fonction $u(S,t)=(H/S)^{2\nu/\sigma^2}C_{\mathrm{BS}}(H^2/S,K,t)$ est solution de l'\omterm{def:dv:black-scholes-three-ways:pde}{équation de Black--Scholes}, ce
que l'on vérifie par substitution. Elle vaut $C_{\mathrm{BS}}(H,K)$ sur $S=H$, et son paiement terminal s'annule pour
$S>H$ lorsque $K\ge H$, car $H^2/S<H\le K$ dans ce cas. La différence est donc la solution qui verse le call au-dessus de la barrière
et zéro sur celle-ci. Dans le log-prix sans dérive, c'est la méthode des images~: la densité des trajectoires tuées en $H$ est la
densité libre moins sa réflexion par rapport à $H$.
\end{proof}

Les sept autres, avec leur parité entrée-sortie $V_{\mathrm{in}}+V_{\mathrm{out}}=V_{\mathrm{vanilla}}$ et leurs rebates,
sont des combinaisons de quatre termes de ce type. La construction implémente les huit et teste la parité pour chaque cas.

\begin{definition}[Rebate de barrière]\label{def:dv:barriers-and-digitals:rebate}
Un \emph{rebate de barrière}\index{rebate de barrière} est un montant fixe versé au détenteur d'une option knock-out lorsque la
barrière est touchée, ou au détenteur d'une option knock-in à l'échéance si elle ne s'est jamais activée.
\end{definition}

\begin{definition}[Option one-touch]\label{def:dv:barriers-and-digitals:onetouch}
Une \emph{option one-touch}\index{option one-touch} verse un montant fixe si le sous-jacent touche la barrière avant
l'échéance, le paiement ayant lieu soit au moment du contact, soit à l'échéance.
\end{definition}

\begin{definition}[Option no-touch]\label{def:dv:barriers-and-digitals:notouch}
Une \emph{option no-touch}\index{option no-touch} verse un montant fixe à l'échéance si le sous-jacent ne touche jamais la
barrière. Une one-touch payée à l'échéance et une no-touch sur la même barrière s'additionnent pour former une obligation zéro-coupon.
\end{definition}

\begin{definition}[Option à double barrière]\label{def:dv:barriers-and-digitals:double}
Une \emph{option à double barrière}\index{option à double barrière} possède une barrière basse et une barrière haute, et s'active ou se désactive
lorsque l'une ou l'autre est touchée. La double no-touch, qui verse un montant si aucune des deux n'est touchée, en est la forme la plus simple.
\end{definition}

Les options touch sont des digitales sur les extrema de la trajectoire. La probabilité de toucher une barrière basse $H<S$ avant $T$ est
$\Phi\bigl((x-\nu T)/\sigma\sqrt T\bigr)+e^{2\nu x/\sigma^2}\Phi\bigl((x+\nu T)/\sigma\sqrt T\bigr)$, avec
$x=\ln(H/S)$. Une double barrière n'admet pas de formule fermée. Sa densité tuée est une série de sinus en log-prix, dont
les termes décroissent comme $e^{-k^2\pi^2\sigma^2T/(2a^2)}$ pour un couloir de largeur logarithmique $a$, si bien que quelques dizaines de termes suffisent.

\begin{example}[Une fiche de barrières]\label{ex:dv:barriers-and-digitals:sheet}
Spot et strike 100, un an, $r=3\%$, $q=1\%$, volatilité 20~\%. Le call vaut 8.83. Avec une option knock-out à 90, elle
vaut 7.23, et la down-and-in 1.60, la somme retombant bien sur le call. Un rebate de 2 versé au contact porte la knock-out à 8.41, puisque
la one-touch à 90 payée au contact vaut 0.592 (0.581 payée à l'échéance). La no-touch à 90 vaut 0.390, la up-and-out
call à barrière 130 vaut 3.10, et la double no-touch entre 85 et 115 vaut 0.141
(la \cref{fig:dv:barriers-and-digitals:barriers}, à gauche).
\end{example}

\section{Couverture statique}

Le delta d'une option à barrière est discontinu à la barrière, comme celui de la digitale à son strike. L'alternative à la
couverture dynamique est un portefeuille d'options vanille qui égale la valeur de l'option à barrière sur la barrière et son paiement à l'échéance, de sorte
qu'elle puisse être dénouée sans coût lorsque la barrière est touchée.

\begin{definition}[Couverture statique]\label{def:dv:barriers-and-digitals:static}
La \emph{couverture statique}\index{couverture statique} réplique une option exotique par un portefeuille d'options vanille figé dès
la mise en place, et réajusté seulement lors d'événements tels qu'une touche de barrière ou l'échéance, plutôt que rééquilibré en continu.
\end{definition}

La symétrie put--call de Carr, Ellis et Gupta fournit les \omterm{def:dv:barriers-and-digitals:static}{couvertures statiques} les plus simples. Avec un carry nul et sans smile, un call
de strike $K$ vaut, sur la barrière $H$, exactement autant que $K/H$ puts de strike $H^2/K$, à chaque instant. Un
call down-and-in est donc couvert en détenant $K/H$ puts de strike $H^2/K$. Si la barrière est touchée, le desk les échange contre
le call sans coût. Si elle ne l'est pas, les puts, dont le strike est inférieur à la barrière, expirent sans valeur, comme l'option knock-in le fait.
Avec un strike de 100, une barrière à 90, un an et une volatilité de 20~\%, la couverture consiste en 1.11 puts de strike 81, valant 1.498, soit exactement
le prix en forme fermée.

La symétrie suppose un smile symétrique, et le smile d'un indice ne l'est pas. Supposons la barrière touchée avec six mois
restants, et le smile à six mois du chapitre 9 conserve sa forme en moneyness. Les 1.11 puts de strike 81 valent alors 2.08 et
le call de strike 100 ne vaut que 0.83. Le skew de volatilité côté baisse valorise les puts, situés sous la barrière, à des
volatilités plus élevées que celle du call situé au-dessus. Une \omterm{def:dv:barriers-and-digitals:static}{couverture statique} construite sur la symétrie aurait coûté bien plus cher que l'option.
Le skew rend caduque toute \omterm{def:dv:barriers-and-digitals:static}{couverture statique} de ce type, indépendante de tout modèle. La couverture calendaire de Derman, Ergener et Kani détient plutôt
des options de plusieurs échéances, choisies à l'aide d'un modèle, de sorte que le portefeuille vaille zéro sur la barrière à un ensemble de
dates donnés.

\begin{example}[Couverture calendaire d'un call up-and-out]\label{ex:dv:barriers-and-digitals:calendar}
Strike 100, barrière 120, un an, 20~\%, carry nul~: le call up-and-out vaut 1.105. Le paiement terminal est un call
au strike 100 moins un call au strike 120 moins 20 digitales au strike 120. En remontant sur $n$ dates régulièrement espacées, le desk ajoute des calls
de strike 120 de chaque échéance, en quantités qui annulent le portefeuille sur la barrière à chaque date. Avec 4 dates, la couverture
coûte 1.77 et vaut jusqu'à 5.4 sur la barrière entre les dates, sur les neuf premiers mois. Avec 16 dates, elle coûte 1.28
et vaut au plus 0.16. Avec 64, elle coûte 1.15 (la \cref{fig:dv:barriers-and-digitals:static}). L'erreur est divisée par deux chaque fois
que le nombre de dates double. Le dernier intervalle avant l'échéance est le plus difficile, car le paiement sur la barrière y saute de
20 à 0.
\end{example}

\begin{omfigure}
\begin{tikzpicture}
\begin{groupplot}[group style={group size=2 by 1,horizontal sep=1.6cm},
  width=6.6cm,height=5cm,
  tick label style={font=\small},label style={font=\small},grid=major,grid style={gray!25},table/col sep=comma]
\nextgroupplot[xlabel={temps (années)},ylabel={valeur de la couverture sur la barrière},xmin=0,xmax=1,ymin=-2,ymax=12,
  title={\small couverture calendaire sur la barrière},
  legend columns=2,legend style={font=\footnotesize,at={(0.5,-0.32)},anchor=north,draw=none,fill=none}]
\addplot[omDef,thick,dashed] table[x=t,y=n4]{figdata/derivatives/15-barriers-and-digitals/calendar_barrier.csv};
\addplot[omThm,very thick] table[x=t,y=n16]{figdata/derivatives/15-barriers-and-digitals/calendar_barrier.csv};
\legend{suivi à 4 dates,suivi à 16 dates}
\nextgroupplot[xmode=log,log basis x=2,xlabel={dates de couverture $n$},ylabel={coût de la couverture},xmin=1.5,xmax=90,ymin=1,ymax=2.5,
  xtick={2,4,8,16,32,64},xticklabels={2,4,8,16,32,64},
  title={\small coût face à la forme fermée},
  legend columns=1,legend cell align=left,legend style={font=\footnotesize,at={(0.5,-0.32)},anchor=north,draw=none,fill=none}]
\addplot[omThm,very thick,mark=*,mark size=1.6pt] table[x=n,y=price]{figdata/derivatives/15-barriers-and-digitals/calendar_prices.csv};
\addplot[gray,thick,dashed] table[x=n,y=uoc]{figdata/derivatives/15-barriers-and-digitals/calendar_prices.csv};
\legend{couverture statique,call up-and-out (1.105)}
\end{groupplot}
\end{tikzpicture}

\omcaption{\omterm{def:dv:barriers-and-digitals:static}{Couverture statique} d'un call up-and-out (strike 100, barrière 120, un an, 20~\%) à partir de calls de plusieurs
échéances. À gauche~: la valeur de la couverture sur la barrière au fil du temps, nulle aux dates sur lesquelles elle est construite, et grande dans le dernier
intervalle, où le paiement saute. À droite~: son coût converge vers le prix de l'option à mesure que les dates se multiplient. Données~: le
tutoriel.}\label{fig:dv:barriers-and-digitals:static}
\end{omfigure}

\section{Surveillance discrète et décalage de barrière}

La plupart des contrats à barrière hors devises sont suivis à des dates discrètes~: clôtures quotidiennes ou fixings. Une option knock-out suivie de
façon discrète survit aux trajectoires qui ont traversé la barrière entre deux observations, si bien qu'elle vaut plus que ne l'annonce la
formule continue. Broadie, Glasserman et Kou ont montré que la formule continue redevient exacte si
la barrière est écartée du spot.

\begin{definition}[Décalage de barrière]\label{def:dv:barriers-and-digitals:shift}
Le \emph{décalage de barrière}\index{décalage de barrière} valorise une option suivie tous les $\Delta t$ avec la
formule continue à la barrière $H\,e^{\pm\beta\sigma\sqrt{\Delta t}}$, écartée du spot, avec $\beta\approx0.5826$.
\end{definition}

La constante vaut $-\zeta(\frac12)/\sqrt{2\pi}$, soit le dépassement moyen attendu d'une marche aléatoire gaussienne au-delà d'un niveau, exprimé en unités
de son pas. La construction la conserve dans une seule constante nommée.

\begin{example}[Précision du décalage]\label{ex:dv:barriers-and-digitals:shift}
Le call knock-out de l'\cref{ex:dv:barriers-and-digitals:sheet} (barrière 90) vaut 7.23 avec un suivi continu. Avec
un suivi quotidien, une simulation de Monte-Carlo de 200\,000 trajectoires donne 7.48 (erreur type 0.03), et la formule décalée 7.45. Avec
un suivi hebdomadaire, elles donnent 7.62 et 7.68, avec un suivi mensuel 8.01 et 8.04, et avec quatre dates 8.33 et 8.38
(la \cref{fig:dv:barriers-and-digitals:barriers}, à droite). Pour un call up-and-out de barrière 120, où le paiement est
élevé à la barrière, le décalage est bon en suivi quotidien (1.26 contre 1.28) et médiocre en suivi trimestriel (2.19 contre 2.50). La correction est un développement en
$\sqrt{\Delta t}$ et échoue lorsque la barrière est à quelques pas du spot ou que le paiement à la barrière est
élevé.
\end{example}

\begin{omfigure}
\begin{tikzpicture}
\begin{groupplot}[group style={group size=2 by 1,horizontal sep=1.6cm},
  width=6.6cm,height=5cm,
  tick label style={font=\small},label style={font=\small},grid=major,grid style={gray!25},table/col sep=comma]
\nextgroupplot[xlabel={spot},ylabel={prix},xmin=88,xmax=130,ymin=0,ymax=32,
  title={\small barrière 90, strike 100},
  legend columns=3,legend style={font=\footnotesize,at={(0.5,-0.32)},anchor=north,draw=none,fill=none}]
\addplot[gray,thick,dashed] table[x=s,y=call]{figdata/derivatives/15-barriers-and-digitals/barrier_curves.csv};
\addplot[omThm,very thick] table[x=s,y=doc]{figdata/derivatives/15-barriers-and-digitals/barrier_curves.csv};
\addplot[omDef,thick] table[x=s,y=dic]{figdata/derivatives/15-barriers-and-digitals/barrier_curves.csv};
\legend{option d'achat,down-and-out,down-and-in}
\nextgroupplot[xmode=log,xlabel={dates de suivi par an},ylabel={call down-and-out},xmin=3,xmax=320,ymin=7,ymax=8.6,
  xtick={4,12,52,252},xticklabels={4,12,52,252},
  title={\small suivi discret},
  legend columns=2,legend style={font=\footnotesize,at={(0.5,-0.32)},anchor=north,draw=none,fill=none}]
\addplot[black,only marks,mark=*,mark size=1.8pt] table[x=n,y=mc]{figdata/derivatives/15-barriers-and-digitals/monitoring.csv};
\addplot[omThm,very thick] table[x=n,y=shifted]{figdata/derivatives/15-barriers-and-digitals/monitoring.csv};
\addplot[gray,thick,dashed] table[x=n,y=cont]{figdata/derivatives/15-barriers-and-digitals/monitoring.csv};
\legend{Monte Carlo,barrière décalée,continu}
\end{groupplot}
\end{tikzpicture}

\omcaption{À gauche~: un call, son down-and-out et son down-and-in selon le spot (un an, 20~\%)~; leur somme redonne le call. À droite~: le
call down-and-out suivi 4, 12, 52 et 252 fois par an~: Monte Carlo contre la formule en temps continu et la formule
à la barrière décalée. Données~: le tutoriel.}\label{fig:dv:barriers-and-digitals:barriers}
\end{omfigure}

La formule à volatilité plate commet une erreur plus grande que le suivi discret lorsque le smile est prononcé. Sous la \omterm{def:dv:local-volatility:model}{volatilité
locale} du chapitre 9, le même knock-out (taux nuls, un an, suivi quotidien) vaut 6.15 par simulation, contre 6.51 selon la
\omterm{thm:dv:black-scholes-three-ways:formula}{formule de Black--Scholes} décalée, à la volatilité à la monnaie de 19.2~\%. La \omterm{def:dv:local-volatility:model}{volatilité locale} est plus élevée aux spots bas, donc
la barrière est touchée plus souvent. L'effet du smile sur ces prix dépend de sa dynamique autant que de son niveau,
\omterm{def:dv:local-volatility:model}{volatilité locale} contre volatilité stochastique (chapitres 9 et 10). Le chapitre 20 évalue les barrières de change avec la
correction propre au marché, construite à partir du \omterm{def:dv:greeks-and-the-hedging-pnl:greeks}{vanna} et du \omterm{def:dv:greeks-and-the-hedging-pnl:greeks}{volga} de l'option.

\section{Risque de gap à la barrière}

Toute couverture de ce chapitre suppose que le prix franchit la barrière~: la \omterm{def:dv:barriers-and-digitals:static}{couverture statique} est dénouée à la barrière,
la couverture dynamique se négocie à proximité, et l'ordre stop-loss est exécuté à son niveau. Un prix qui saute par-dessus la
barrière met en échec les trois.

\begin{definition}[Risque de gap]\label{def:dv:barriers-and-digitals:gap}
Le \emph{risque de gap}\index{risque de gap} est le risque qu'un prix se déplace de manière discontinue à travers un niveau auquel la valeur d'une position
ou de sa couverture change, une barrière, un strike proche de l'échéance ou un stop, de sorte que la couverture prévue à ce niveau soit exécutée à un
moins bon prix, voire pas du tout.
\end{definition}

Le knock-out inversé, dont le paiement est élevé à la barrière, porte ce risque sous sa forme la plus pure.

\begin{example}[Un gap à travers une barrière inversée]\label{ex:dv:barriers-and-digitals:gap}
Un call up-and-out, strike 100, barrière 120, a encore un mois à courir, et le spot vaut 119. Il vaut 1.38 et son delta est
$-1.37$~: une hausse rapproche la barrière et détruit le paiement
(la \cref{fig:dv:barriers-and-digitals:gap}). Le vendeur se couvre en vendant 1.37 actions. Si le spot passe à 119.5, le
P\&L de la position couverte vaut $+0.005$. S'il saute à 125, l'option disparaît, et le vendeur encaisse ses 1.38, mais les
actions vendues à découvert perdent $1.37\times6=8.19$~: une perte nette de 6.81, soit cinq fois la prime.
\end{example}

\begin{omfigure}
\begin{tikzpicture}
\begin{groupplot}[group style={group size=2 by 1,horizontal sep=1.6cm},
  width=6.6cm,height=5cm,
  tick label style={font=\small},label style={font=\small},grid=major,grid style={gray!25},table/col sep=comma]
\nextgroupplot[xlabel={spot},ylabel={valeur},xmin=95,xmax=121,ymin=0,ymax=9,
  title={\small call up-and-out, un mois}]
\addplot[omThm,very thick] table[x=s,y=value]{figdata/derivatives/15-barriers-and-digitals/reverse.csv};
\draw[gray,dashed] (axis cs:120,0) -- (axis cs:120,9);
\nextgroupplot[xlabel={spot},ylabel={delta},xmin=95,xmax=121,ymin=-1.6,ymax=0.8,
  title={\small son delta}]
\addplot[omDef,very thick] table[x=s,y=delta]{figdata/derivatives/15-barriers-and-digitals/reverse.csv};
\draw[gray,dashed] (axis cs:120,-1.6) -- (axis cs:120,0.8);
\draw[gray] (axis cs:95,0) -- (axis cs:121,0);
\end{groupplot}
\end{tikzpicture}

\omcaption{Un knock-out inversé près de sa barrière~: strike 100, barrière 120 (en tirets), un mois à courir, 20~\%. Sa valeur
culmine bien avant la barrière et son delta devient fortement négatif~; un hedger vendeur de l'option est vendeur du sous-jacent
là, et perd sur tout gap à travers la barrière. Données~: le code du chapitre.}\label{fig:dv:barriers-and-digitals:gap}
\end{omfigure}

Les gaps ne tiennent pas qu'aux week-ends et aux annonces de résultats. Le 15 janvier 2015, la Banque nationale suisse a abandonné son
cours plancher de 1.20 franc suisse pour un euro. Le taux de référence de la BCE pour le franc valait 1.2010 le 14 et
1.0280 le 15, soit une chute de 14.4~\% entre deux fixings, et 1.0128 le 16. Une barrière située n'importe où dans cette
fourchette a été franchie d'un seul mouvement, et une couverture prévue à la barrière ne pouvait s'exécuter qu'au prochain cours échangé. Un
système de plancher que défend une banque centrale ressemble à la barrière la plus sûre du marché, et, quand il cède, il cède par un gap. Les desks évaluent le risque avec des modèles à sauts
(chapitre 13) et constituent des réserves pour ce risque (chapitre 27).

\section{Tutoriel~: barrières et options digitales}\label{tut:dv:barriers-and-digitals:tut}

\begin{tutorial}
\textbf{Objectif.} Évaluer une \omterm{def:dv:barriers-and-digitals:digital}{option digitale} avec et sans smile, ainsi que sa \omterm{def:dv:barriers-and-digitals:overhedge}{surcouverture par call spread}~; évaluer les barrières sous forme fermée,
sous suivi discret et sous \omterm{def:dv:local-volatility:model}{volatilité locale}~; construire deux \omterm{def:dv:barriers-and-digitals:static}{couvertures statiques}~; mesurer le gap. \textbf{État final~:} les
quatre figures et les chiffres du problème du week-end.
\begin{tutsteps}
\item \textbf{Une \omterm{def:dv:barriers-and-digitals:digital}{option digitale} sous smile} est l'opposée de la dérivée du call par rapport au strike~:
\omcode{code/firm/barrier/firm_barrier.py}{34}{38}{Option digitale cohérente avec le smile.}\label{lst:dv:barriers-and-digitals:digital}
\item \textbf{Le \omterm{def:dv:barriers-and-digitals:shift}{décalage de barrière} et sa vérification.} La barrière décalée, et une simulation de Monte-Carlo de barrières
à suivi discret, avec des pas lognormaux exacts~:
\omcode{code/firm/barrier/firm_barrier.py}{129}{150}{Décalage de barrière et Monte-Carlo à suivi discret.}\label{lst:dv:barriers-and-digitals:shift}
\item \textbf{Une \omterm{def:dv:barriers-and-digitals:notouch}{option no-touch} double} par sa série en sinus~:
\omcode{code/firm/barrier/firm_barrier.py}{110}{125}{Option no-touch double.}\label{lst:dv:barriers-and-digitals:dnt}
\item \textbf{Exécuter} \texttt{dv\_barrier.digital\_prices()}, \texttt{overhedge(200)}, \texttt{monitoring()},
\texttt{local\_vol\_doc()}, \texttt{calendar\_example()}, \texttt{reverse\_barrier()} et \texttt{fig\_barrier.py}.
\end{tutsteps}
\textbf{Pour aller plus loin.} Évaluer l'\omterm{def:dv:barriers-and-digitals:digital}{option digitale} comme un call spread centré sur le strike et comparer son risque avec la
surcouverture~; remplacer, dans la couverture calendaire, les calls à la barrière par des puts~; déplacer la barrière up-and-out à 110 et relancer
l'analyse du gap.
\end{tutorial}

\section{Construction~: barrières et options digitales}\label{bld:dv:barriers-and-digitals:barrier}

\begin{build}
\bfield{Rôle} L'outil de pricing des barrières et des \omterm{def:dv:barriers-and-digitals:digital}{options digitales} de la firme miniature~: formules fermées avec rebates, touchers, double no-touch,
\omterm{def:dv:barriers-and-digitals:shift}{décalage de barrière} pour surveillance discrète et un constructeur de \omterm{def:dv:barriers-and-digitals:static}{couverture statique}~; il sert de socle aux autocallables du chapitre 18 et aux
barrières de change du chapitre 20.
\bfield{Interface} \texttt{digital}, \texttt{digital\_delta}, \texttt{digital\_smile(s, k, t, r, q, vol\_of\_k)}~;
\texttt{call\_spread(\ldots, width, notional, below)}~; \texttt{barrier(s, k, h, t, r, q, vol, kind, right, rebate)} avec
\texttt{kind} prenant les valeurs down-out, down-in, up-out, up-in~; \texttt{hit\_probability}, \texttt{one\_touch(\ldots, at\_hit)},
\texttt{no\_touch}, \texttt{double\_no\_touch}~; \texttt{bgk\_shift}, \texttt{BGK\_BETA}~; \texttt{mc\_barrier}~;
\texttt{symmetry\_down\_in\_call}~; \texttt{calendar\_hedge}, \texttt{calendar\_hedge\_on\_barrier}.
\bfield{Règles} Les options knock-in et knock-out s'additionnent toujours pour redonner l'option vanille (testé sur les huit cas)~; la surveillance discrète est pricée avec
le décalage ou par simulation, jamais avec la formule continue~; une \omterm{def:dv:barriers-and-digitals:digital}{option digitale} est toujours pricée avec la pente du smile~; une
\omterm{def:dv:barriers-and-digitals:digital}{option digitale} vendue est enregistrée à sa surcouverture.
\bfield{Tests d'acceptation} \texttt{code/firm/barrier/tests/}~: parité in-out pour chaque type, chaque droit et chaque strike~; une
barrière lointaine redonne l'option vanille~; la formule décalée coïncide avec un Monte-Carlo quotidien~; les identités de toucher, et la double no-touch
face à sa limite à barrière simple et à une simulation~; l'\omterm{def:dv:barriers-and-digitals:digital}{option digitale} avec smile se réduit à la version plate sans skew, et les call
spreads encadrent l'\omterm{def:dv:barriers-and-digitals:digital}{option digitale}~; la couverture par symétrie price exactement le down-and-in, et la couverture calendaire converge.
\bfield{Pour aller plus loin} Barrières sous \omterm{def:dv:local-volatility:model}{volatilité locale} et stochastique par différences finies (chapitre 22)~; barrières à fenêtre~;
la correction vanna--volga du chapitre 20.
\end{build}

\begin{omsources}
\item M.~Broadie, P.~Glasserman et S.~Kou, ``A continuity correction for discrete barrier options'', \emph{Mathematical
Finance} 7(4) (1997) 325--349.
\item E.~Derman, D.~Ergener et I.~Kani, ``Static options replication'', \emph{Journal of Derivatives} 2(4) (1995) 78--95.
\item P.~Carr, K.~Ellis et V.~Gupta, ``Static hedging of exotic options'', \emph{Journal of Finance} 53(3) (1998)
1165--1190.
\item Swiss National Bank, communiqué de presse, 15 janvier 2015~; European Central Bank, euro foreign exchange reference rates.
\end{omsources}

\section{Exercices}

\begin{exercise}[$\star$]\label{exo:dv:barriers-and-digitals:1}
Montrer qu'un call digital cash et un put digital cash de même strike et même échéance s'additionnent en un zéro-coupon. Que
cela implique-t-il pour leurs deltas~?
\end{exercise}

\begin{exercise}[$\star$]\label{exo:dv:barriers-and-digitals:2}
Pourquoi un skew à la baisse rend-il un call digital plus cher que son prix de Black--Scholes à la volatilité du strike~?
\end{exercise}

\begin{exercise}[$\star$]\label{exo:dv:barriers-and-digitals:3}
Un call knock-out et son knock-in sont pricés à 7.23 et 1.70 alors que l'option vanille vaut 8.83. Qu'est-ce qui ne va pas~?
\end{exercise}

\begin{exercise}[$\star\star$]\label{exo:dv:barriers-and-digitals:4}
À carry nul, vérifier numériquement la couverture par symétrie~: pricer 1.11 puts de strike 81 et le call down-and-in de
barrière 90 à 20~\% de volatilité. Pourquoi les puts expirent-ils sans valeur chaque fois que la barrière n'est pas touchée~?
\end{exercise}

\begin{exercise}[$\star\star$]\label{exo:dv:barriers-and-digitals:5}
Une barrière est surveillée quotidiennement à 20~\% de volatilité. De combien le décalage déplace-t-il une barrière à 90~? À 120~?
\end{exercise}

\begin{exercise}[$\star\star$]\label{exo:dv:barriers-and-digitals:6}
Expliquer pourquoi le call up-and-out de l'\cref{ex:dv:barriers-and-digitals:gap} a un delta négatif au voisinage de la barrière, et pourquoi
sa valeur culmine sous la barrière.
\end{exercise}

\begin{exercise}[$\star\star\star$]\label{exo:dv:barriers-and-digitals:7}
\emph{Codage.} Relancer la comparaison en surveillance discrète pour un put down-and-out de barrière 90. Le décalage fonctionne-t-il aussi bien
que pour le call~?
\end{exercise}

\begin{exercise}[$\star\star\star$]\label{exo:dv:barriers-and-digitals:8}
\emph{Trouver le défaut.} «~Notre livre de barrières est sûr~: chaque barrière a un ordre stop-loss en attente à son niveau, de sorte que la
couverture est exécutée exactement au moment où l'option se déclenche.~»
\end{exercise}

\section{Problème~: l'option digitale à l'échéance}

\begin{problem}[{Problème du week-end --- quelle largeur pour un call spread}]\label{pb:dv:barriers-and-digitals:1}
Un client achète une \omterm{def:dv:barriers-and-digitals:digital}{option digitale} à trois mois qui verse 10 millions si l'indice, aujourd'hui à 5\,000, clôture au-dessus de 5\,000 à l'échéance. Le
desk pricera sur la surface du chapitre 9 mise à l'échelle de 5\,000 et enregistrera l'\omterm{def:dv:barriers-and-digitals:digital}{option digitale} au prix d'une \omterm{def:dv:barriers-and-digitals:overhedge}{surcouverture par call spread}. Sa
limite~: une erreur d'un point sur le prix de règlement ne doit pas faire bouger la couverture de plus de 50\,000.

\medskip\noindent\textbf{Partie I --- L'\omterm{def:dv:barriers-and-digitals:digital}{option digitale}.}
\begin{enumerate}
\item Donner le prix de l'\omterm{def:dv:barriers-and-digitals:digital}{option digitale} avec la volatilité à la monnaie et avec le smile.
\item Quel est le terme de \omterm{def:dv:greeks-and-the-hedging-pnl:greeks}{vega}, et pourquoi est-il positif~?
\item Une minute avant l'échéance, avec l'indice à 4\,999.9 et une volatilité de 20~\%, donner le delta de la digitale par point et en
notionnel d'indice.
\item Pourquoi ce delta ne peut-il pas être négocié~?
\item Quel est le gamma de la digitale juste en dessous et juste au-dessus du strike à proximité de l'échéance~?
\end{enumerate}

\medskip\noindent\textbf{Partie II --- La surcouverture.}
\begin{enumerate}[resume]
\item Quel call spread domine la digitale, et pourquoi est-ce celui que le vendeur achète~?
\item Quelle largeur respecte la limite du desk~?
\item Donner le prix du spread et la surcouverture facturée au client.
\item Donner la surcouverture pour des largeurs de 50, 100 et 400 points.
\item Que reçoit le client, par rapport à la digitale, si l'indice se règle à 4\,900~?
\end{enumerate}

\medskip\noindent\textbf{Partie III --- Barrières.}
\begin{enumerate}[resume]
\item Donner le call down-and-out à un an (strike 100, barrière 90, 20~\%, $r=3\%$, $q=1\%$) et son knock-in.
\item Le donner avec surveillance quotidienne, par simulation et par décalage.
\item Le donner sous la \omterm{def:dv:local-volatility:model}{volatilité locale} du chapitre 9 et expliquer l'écart avec le prix à volatilité plate.
\item Donner la double no-touch entre 85 et 115.
\item Pricer la one-touch à 90 payée au touché et à l'échéance, et expliquer la différence.
\end{enumerate}

\medskip\noindent\textbf{Partie IV --- Jugement.}
\begin{enumerate}[resume]
\item Pourquoi la surcouverture est-elle équitable pour le client~?
\item Quand utiliser plutôt un call spread centré~?
\item Comment couvrir un reverse knock-out à proximité de sa barrière~?
\item Énoncer le \emph{résultat nommé}~: la largeur du call spread qui plafonne la perte du hedger à 50\,000 par point de règlement sur
la digitale de 10 millions, et la surcouverture qu'elle facture au client.
\item En une phrase~: qu'ont en commun les digitales et les barrières~?
\end{enumerate}
\end{problem}

\section{Questions d'entretien}

\begin{interviewq}[$\star$ \iqroles{trader}]\label{iq:dv:barriers-and-digitals:1}
Comment pricez-vous et couvrez-vous une \omterm{def:dv:barriers-and-digitals:digital}{option digitale}~?
\end{interviewq}

\begin{interviewq}[$\star$ \iqroles{chercheur}]\label{iq:dv:barriers-and-digitals:2}
Pouvez-vous dériver le prix d'un call down-and-out avec le principe de réflexion~?
\end{interviewq}

\begin{interviewq}[$\star\star$ \iqroles{trader, chercheur}]\label{iq:dv:barriers-and-digitals:3}
Pourquoi le skew compte-t-il pour une \omterm{def:dv:barriers-and-digitals:digital}{option digitale}, et dans quel sens~?
\end{interviewq}

\begin{interviewq}[$\star\star$ \iqroles{développeur}]\label{iq:dv:barriers-and-digitals:4}
Votre Monte-Carlo price un knock-out à surveillance quotidienne au-dessus de la forme fermée. Est-ce un bug~?
\end{interviewq}

\begin{interviewq}[$\star\star$ \iqroles{risque}]\label{iq:dv:barriers-and-digitals:5}
Qu'est-ce que le \omterm{def:dv:barriers-and-digitals:gap}{risque de gap} sur un livre de barrières, et comment le mesureriez-vous et le limiteriez-vous~?
\end{interviewq}

\begin{interviewq}[$\star\star\star$ \iqroles{trader, chercheur}]\label{iq:dv:barriers-and-digitals:6}
Comment construiriez-vous une \omterm{def:dv:barriers-and-digitals:static}{couverture statique} pour un call down-and-in~? Quand échoue-t-elle~?
\end{interviewq}
